\section{Related Work}
\label{sec:related-work}

The relationship between judging subjectivity and competitive scoring has been
studied in other adjudicated sports. Wolframm~\cite{wolframm2023bias} examined
elite dressage competitions and reported significant associations between scores
and contextual variables including prior ranking, starting order, home status,
and competitor nationality across seven events; the study acknowledges
alternative explanations and sample-size constraints. We cite this only as
evidence that structured non-performance factors have been examined in another
judged sport, not as a direct analogue to indoor percussion.

Trajectory modeling of sports performance across a competitive season offers a
closer methodological parallel. Dolmeta, Argiento, and
Montagna~\cite{dolmeta2023garch} apply Bayesian GARCH models to intra- and
inter-season shot-put trajectories, demonstrating that early-season performance
can inform later-season predictions. The authors explicitly scope their
framework to sports with measurable physical outcomes; SCPA scores are
adjudicated constructs, not physical measurements, which limits direct
methodological transfer.

Within the activity domain, Winter Guard International (WGI) operates color
guard, percussion, and winds competitions~\cite{wgiabout2026}. Its percussion
season culminates in the Percussion World Championships, providing a broader
championship context for ensembles that also compete in SCPA
events~\cite{wgipercussion2026}. SCPA advertises its competitions as using
``WGI format''~\cite{scpaabout2026}; accordingly, the current WGI percussion
caption structure is directly relevant to interpreting SCPA scores. However,
the available sources do not establish that SCPA used unchanged, identical
score sheets in every season from 2017 through 2026.

bandScores~\cite{bandscores2023probabilistic} uses regional score timing to
discuss championship advancement in WGI color guard. That nonacademic article
concerns WGI color guard, whereas the present data contain SCPA percussion
performances only; it is an informal forecasting example, not peer-reviewed
evidence or external validation.

This paper introduces a newly engineered competitive-percussion dataset and
applies debut-to-championship prediction within it. The contribution is the
dataset and the application; the regression and boosting techniques used are
standard and not claimed as novel.
